\documentclass[11pt]{article}
\usepackage[a4paper,margin=1in]{geometry}
\usepackage{amsmath,amssymb,amsthm}
\usepackage{graphicx}
\usepackage{xcolor}
\usepackage{hyperref}
\usepackage{microtype}

\newtheorem{theorem}{Theorem}
\newtheorem{lemma}{Lemma}
\newcommand{\cube}{\{-1,1\}}
\newcommand{\norm}[1]{\left\lVert#1\right\rVert}
\newcommand{\R}{\mathbb R}

\begin{document}

\begin{center}
{\Large\bf Exact homogeneous-level projection norms on Boolean cubes}\par
\vspace{0.7em}
Automated mathematical discovery run \texttt{fa\_banach\_001}\par
\vspace{0.35em}
13 August 2026
\end{center}

\medskip
\noindent\fbox{\parbox{0.94\linewidth}{
\textbf{Status.} Candidate full solution, likely valid, subject to human review.  The result
answers the source's Section~4.3 question affirmatively and proves the stronger exact
formula for every fixed degree and Fourier level.
}}

\section{Source question}

The source is A. Defant, M. Masty\l o, and A. P\'erez, \emph{On the Fourier
spectrum of functions on Boolean cubes}, arXiv:1706.03670~\cite{DMP}.  On
source p.~24, Section~4.3 asks whether $1+\sqrt2$ is the best constant $C$
such that
\[
  \norm{f_m}_\infty\le C^d\norm f_\infty
\]
for every real $f:\cube^n\to\R$ of Fourier degree at most $d$ and every
$0\le m\le d$.

\begin{center}
\includegraphics[width=0.96\linewidth]{figures/open_problem_crop.png}
\end{center}

Write the Fourier--Walsh polynomial of $f$ as
\[
 F(x)=\sum_{S\subseteq[n]}\widehat f(S)x^S,
 \qquad x^S=\prod_{j\in S}x_j,
\]
and let $F_m=f_m=\sum_{|S|=m}\widehat f(S)x^S$.  The polynomial $F$ is
multi-affine, and therefore
\begin{equation}\label{eq:normequality}
 \norm F_{[-1,1]^n}=\norm f_{\cube^n}.
\end{equation}

For $0\le m\le d$, define the classical coefficient constant
\begin{equation}\label{eq:markovdef}
 M_{m,d}:=\sup\left\{|[t^m]p|:\deg p\le d,
       \ \sup_{-1\le t\le1}|p(t)|\le1\right\}.
\end{equation}
Markov's sharp coefficient theorem says that $M_{m,d}$ is the absolute
$t^m$ coefficient of $T_d$ when $m\equiv d\pmod2$, and of $T_{d-1}$
otherwise~\cite[Eqns.~(12)--(13)]{DMP}; see also~\cite[Ch.~3]{BE}.

\section{Proof intuition}

The known upper bound reduces the multivariable question to one variable:
on each radial line, $F(tx)$ is a degree-$d$ polynomial whose $t^m$
coefficient is $F_m(x)$, so Markov's theorem gives the constant $M_{m,d}$.
What was missing is sharpness inside the multi-affine class.

To recover a one-variable Chebyshev extremizer without repeated variables,
spread its single variable over more and more Boolean coordinates:
\[
 f_n(\varepsilon)=T_{d'}\!\left(\frac{\varepsilon_1+\cdots+\varepsilon_n}{n}\right).
\]
This remains bounded by one.  After reducing powers by $x_j^2=1$, the
degree-$m$ part of the $m$th power of the average is asymptotically the
ordinary squarefree term: collisions among indices have lower order in $n$.
All higher Chebyshev powers contribute only collision terms and vanish in the
limit.  The level-$m$ value therefore tends exactly to the extremal Chebyshev
coefficient.

\section{Exact result}

For fixed $m,d$, put
\[
 \Lambda_{m,d}:=sup_{n\ge1}\ \sup_{\substack{0\ne f:\cube^n\to\R\\
        \deg f\le d}}
 \frac{\norm{f_m}_\infty}{\norm f_\infty}.
\]

\begin{theorem}\label{thm:exact}
For every $d\ge1$ and $0\le m\le d$,
\[
 \boxed{\Lambda_{m,d}=M_{m,d}.}
\]
Consequently, the smallest $C$ for which
$\norm{f_m}_\infty\le C^d\norm f_\infty$ holds uniformly in $n,d,m$ is
\[
 \boxed{C=1+\sqrt2.}
\]
\end{theorem}

\begin{proof}
We first prove the upper bound.  Let $F$ be the multi-affine polynomial of a
function $f$ of degree at most $d$.  For $x\in[-1,1]^n$, set
\[
 p_x(t):=F(tx)=\sum_{r=0}^d t^rF_r(x).
\]
If $t\in[-1,1]$, then $tx\in[-1,1]^n$.  Hence
$\norm{p_x}_{[-1,1]}\le\norm F_{[-1,1]^n}$, while the coefficient of $t^m$
in $p_x$ is $F_m(x)$.  The definition~\eqref{eq:markovdef}, followed by
\eqref{eq:normequality}, gives
\[
 |F_m(x)|\le M_{m,d}\norm F_{[-1,1]^n}
              =M_{m,d}\norm f_{\cube^n}.
\]
Taking the supremum over $x$ proves $\Lambda_{m,d}\le M_{m,d}$.

For the reverse inequality, choose
\[
 d'=\begin{cases}d,&m\equiv d\pmod2,\\d-1,&m\not\equiv d\pmod2,
 \end{cases}
 \qquad
 T_{d'}(t)=\sum_{k=0}^{d'}a_kt^k.
\]
Then $|a_m|=M_{m,d}$.  For each $n$, define
\begin{equation}\label{eq:fn}
 f_n(\varepsilon)=T_{d'}\!\left(
           \frac{\varepsilon_1+\cdots+\varepsilon_n}{n}\right),
 \qquad \varepsilon\in\cube^n.
\end{equation}
The argument of $T_{d'}$ lies in $[-1,1]$, so $\norm{f_n}_\infty\le1$.
Let $Q_n$ be the unique multi-affine polynomial agreeing with~\eqref{eq:fn}
on $\cube^n$; equivalently, reduce every $x_j^r$ using $x_j^2=1$.
Thus $Q_n$ is precisely the Fourier--Walsh polynomial of $f_n$ and has degree
at most $d'$.

For integers $k,m\ge0$, let $N_{n,k,m}$ be the number of ordered $k$-tuples
$(i_1,\ldots,i_k)\in[n]^k$ for which exactly $m$ indices occur an odd number
of times.  Expanding $(x_1+\cdots+x_n)^k$ and reducing modulo
$x_j^2-1$ shows that
\begin{equation}\label{eq:countidentity}
 (Q_n)_m(\mathbf1)=\sum_{k=0}^{d'}a_kn^{-k}N_{n,k,m}.
\end{equation}
If $k<m$, then $N_{n,k,m}=0$.  If $k=m$, each of the $m$ occurring indices
must occur once, and hence
\begin{equation}\label{eq:maincount}
 N_{n,m,m}=(n)_m=n(n-1)\cdots(n-m+1).
\end{equation}
Finally suppose $k>m$.  If such a tuple uses $r$ distinct indices, $m$ of
them have odd positive multiplicity and the other $r-m$ have even
multiplicity at least two.  Therefore
\[
 k\ge m+2(r-m),\qquad r\le\frac{k+m}{2}.
\]
For fixed $k,m$, choosing the distinct indices and their multiplicities gives
the uniform estimate
\begin{equation}\label{eq:errorcount}
 N_{n,k,m}=O_{k,m}\!\left(n^{(k+m)/2}\right).
\end{equation}
Equations~\eqref{eq:countidentity}--\eqref{eq:errorcount} now imply
\[
 (Q_n)_m(\mathbf1)
 =a_m\frac{(n)_m}{n^m}
  +\sum_{k=m+1}^{d'}a_kO_{k,m}\!\left(n^{-(k-m)/2}\right)
 \longrightarrow a_m.
\]
Since $\norm{f_n}_\infty\le1$,
\[
 \Lambda_{m,d}\ge\limsup_{n\to\infty}
   \norm{(f_n)_m}_\infty
 \ge\lim_{n\to\infty}|(Q_n)_m(\mathbf1)|
 =|a_m|=M_{m,d}.
\]
This proves the exact identity.

If ``degree $d$'' is read as \emph{exactly} $d$ rather than at most $d$,
the only extra case to note is $d'=d-1$ (where necessarily $m<d$).  On $d$
fresh Boolean coordinates $\eta$, replace $f_n$ by
\[
 f_n(\varepsilon)+\delta\eta_1\cdots\eta_d.
\]
Its degree is exactly $d$, its level-$m$ part is unchanged, and its norm is at
most $1+\delta$.  Letting $\delta\downarrow0$ gives the same lower bound.

It remains to identify the best common exponential base.  Put
$A=1+\sqrt2$, so $A^{-1}=\sqrt2-1$.  The nonzero coefficients of $T_r$
alternate in sign along their common parity.  Consequently all terms have
the same phase at $i$, and the standard closed form for $T_r$ gives
\begin{equation}\label{eq:Ti}
 \sum_{j=0}^r|[t^j]T_r|
 =|T_r(i)|
 =\frac{A^r+(-1)^rA^{-r}}2.
\end{equation}
In particular every coefficient of $T_r$ is at most $A^r$.  The same bound
for $T_d$ and $T_{d-1}$ gives $M_{m,d}\le A^d$, so $C=A$ works.
Conversely,~\eqref{eq:Ti} and the fact that $T_r$ has at most $r+1$
coefficients show that some coefficient has magnitude at least
$(A^r-A^{-r})/(2(r+1))$.  Its $r$th root tends to $A$.  Thus no smaller
uniform base can work.  This completes the proof.
\end{proof}

\section{Checks, scope, and novelty audit}

The proof is exact and has no computational dependency.  As a sanity check,
the symmetric Fourier formula was evaluated for $(d,m)=(6,2),(7,3),(10,6)$;
the level values converge to the predicted Chebyshev coefficients
$18,56,1120$.  The script in \texttt{code/check\_chebyshev\_limits.py}
reproduces this check.

The result answers only Section~4.3.  It does not resolve the distinct
Aaronson--Ambainis conjecture or the source's question about subexponential
Lorentz Bohnenblust--Hille constants.

A bounded search on 13 August 2026 covered the arXiv identifier and title,
the exact Section~4.3 wording, and combinations of ``homogeneous parts'',
``Boolean cube'', ``Fourier level projection'', ``Markov coefficient'', and
``Chebyshev''.  It found later general dimension-free bounded-projection work,
including~\cite{BKSZ}, but no exact statement of
$\Lambda_{m,d}=M_{m,d}$ and no explicit resolution of Section~4.3.
Novelty confidence is therefore provisional.  Human review should focus on
the Fourier/multilinearization identification and the count
in~\eqref{eq:errorcount}.

\begin{thebibliography}{9}
\bibitem{DMP}
A. Defant, M. Masty\l o, and A. P\'erez,
\emph{On the Fourier spectrum of functions on Boolean cubes},
Math. Ann. 374 (2019), 653--680; arXiv:1706.03670.

\bibitem{BE}
P. Borwein and T. Erd\'elyi,
\emph{Polynomials and Polynomial Inequalities},
Graduate Texts in Mathematics 161, Springer, 1995.

\bibitem{BKSZ}
L. Becker, O. Klein, J. Slote, A. Volberg, and H. Zhang,
\emph{Dimension-free discretizations of the uniform norm by small product sets},
Invent. Math. 239 (2025), 469--503; arXiv:2310.07926.
\end{thebibliography}

\end{document}
